\section{Literal local periods and actual marked quotients}
A period appearing among global relators need not be the least period
of the corresponding parabolic. This distinction is checked explicitly.
For a finite local group generated by $u,v$, take one of the three
standard rotation presentations
\begin{align*}
 T_{63}(a,b)&=\gen{u,v\mid u^6,v^3,(u^{-1}v)^2,
                (v^{-1}u^{-2})^a(vu^2)^b},\\
 T_{36}(a,b)&=\gen{u,v\mid u^3,v^6,(u^{-1}v)^2,
                (u^{-1}v^{-2})^a(uv^2)^b},\\
 T_{333}(a,b)&=\gen{u,v\mid u^3,v^3,(u^{-1}v)^3,
                (uv^{-1}u)^a(uv)^b}.
\end{align*}
These retain the exact source parameter conventions. Their standard orders
are $6(a^2+ab+b^2)$, $6(a^2+ab+b^2)$ and $3(a^2+ab+b^2)$.
The usual faithful Euclidean triangle realization is retained here;
for each finite row the program also computes the abstract local order
independently and tests the exact marked epimorphism into the parabolic.
Consequently no local isomorphism is inferred from an unmarked order match.

Let $U,V$ denote the two translation words in the last relator of the
chosen presentation. The program tests that their actual images commute,
generate a normal subgroup $L$, and have index six (or three) in the
local group. It determines $m=|U|$ and the least $b>0$ for which
$U^aV^b=1$ with $0\leq a<m$. Then
\[
 \ker\bigl(\Z^2\longrightarrow L,\ (i,j)\longmapsto U^iV^j\bigr)
       =\Z(m,0)+\Z(a,b).
\]
Indeed reduce the second coordinate modulo $b$ using $(a,b)$ and then
the first modulo $m$. Minimality leaves exactly $mb$ different images;
the program independently verifies $mb=|L|$. This proves the entire
lattice kernel, not merely two relations. The order of the complete
abstract local presentation after adjoining $U^m$ and $U^aV^b$ is then
computed and compared to the actual parabolic order. The generator map
is onto and the equal orders prove an exact marked isomorphism.

For spherical rank-three residues the same method uses the complete
triangle presentation $\gen{u,v\mid u^p,v^q,(u^{-1}v)^r}$, with
$1/p+1/q+1/r>1$, and an explicit onto generator map. Its exact order is
computed. Every local result therefore specifies its marking and not
just a name or a number. The full coset tests in the preceding section
independently ensure that these are the actual residues of the global
geometry.

When a printed word does not impose its claimed torus period, it is
kept literally in the global presentation. The certificate computes
the actual local lattice from the three triangle relations and the
actual finite parabolic. It never inserts an intended torus word
silently. Such a case is recorded separately from a proper collapse
of an otherwise valid local period presentation.
